% 核心结论：有限跳跃在导数中产生与跳跃量成正比的界面分布。
% 图契约：物理竞赛讲义；印刷宽16cm，最小字号8.5pt；几何示意与公式分区。
% 布局：[几何/曲线 7cm] [定义/结果 7cm]；实线主量，虚线辅助；无装饰数据。
\documentclass[tikz,border=6pt]{standalone}
\usepackage{fontspec,xeCJK,amsmath,amssymb,graphicx}
\IfFontExistsTF{Times New Roman}{\setmainfont{Times New Roman}}{\setmainfont{TeX Gyre Termes}}
\IfFontExistsTF{Heiti SC}{\setCJKmainfont[AutoFakeBold=1.5]{Heiti SC}}{\setCJKmainfont[AutoFakeBold=1.5]{FandolHei-Regular}}
\usetikzlibrary{arrows.meta,bending,calc,positioning,patterns}
\definecolor{ink}{HTML}{243745}
\definecolor{blue}{HTML}{275DAB}
\definecolor{green}{HTML}{177C63}
\definecolor{red}{HTML}{BD4935}
\definecolor{orange}{HTML}{AA6C16}
\definecolor{purple}{HTML}{715393}
\tikzset{
 every node/.style={font=\fontsize{9.5}{13}\selectfont,text=ink},
 axis/.style={-{Stealth[length=4pt]},line width=.65pt,draw=ink!65},
 vec/.style={-{Stealth[length=5pt]},line width=1pt,draw=blue},
 short/.style={-{Stealth[length=3pt]},line width=.8pt},
 guide/.style={dashed,draw=ink!40,line width=.5pt},
 card/.style={rounded corners=2pt,draw=blue!60,fill=blue!4,
   align=center,text width=6.8cm,inner xsep=9pt,inner ysep=7pt},
 title/.style={font=\fontsize{11}{14}\selectfont\bfseries,text width=15cm,align=center},
 label/.style={font=\fontsize{8.5}{11}\selectfont,align=center}}
\begin{document}
\begin{tikzpicture}[x=1cm,y=1cm]
\path[use as bounding box] (-8,-3.35) rectangle (8,3.35);
\node[title] at (0,2.95) {阶跃的分布导数与界面跳跃};
\begin{scope}[shift={(-4,0)}]

\draw[axis] (-2.5,-1.2)--(2.6,-1.2) node[right] {$x$};\draw[axis] (-2.2,-1.5)--(-2.2,2) node[above] {$H$};
\draw[blue,line width=1pt] (-2,-1.2)--(0,-1.2);\draw[blue,line width=1pt] (0,.8)--(2.2,.8);
\draw[guide] (0,-1.2)--(0,.8);\node[below] at (0,-1.2) {$a$};
\draw[vec,red] (0,-1.1)--(0,1.8);\node[label,red] at (1.4,1.7) {$\delta(x-a)$\\箭头表示单位权重};

\end{scope}\node[card,draw=blue!65,fill=blue!4] at (3.9,1.65) {$\displaystyle \frac{d}{dx}H(x-a)=\delta(x-a)$};
\node[card,draw=green!65,fill=green!4] at (3.9,0) {$F=F^-+[F]H(x-a)$\\$F^\prime=(F^\prime)_{\rm reg}+[F]\delta(x-a)$};
\node[card,draw=orange!65,fill=orange!4] at (3.9,-1.65) {曲面界面：\\$\nabla F=(\nabla F)_{\rm reg}+[F]\hat{\mathbf n}\,\delta_S$};
\end{tikzpicture}\end{document}
